% Clase 25 --- Onda esferica: la intensidad cae como 1/r^2 (§5).
% El argumento es conservacion de la energia: la MISMA potencia atraviesa los dos
% cascarones, pero repartida en areas que crecen como r^2. El dibujo tiene que
% mostrar el cono de energia ensanchandose, no solo las esferas.
\begin{center}
\begin{tikzpicture}[scale=1]

  % fuente
  \filldraw[figamber] (0,0) circle (2.4pt);
  \node[figlbl, figamber, anchor=north east] at (-0.2,-0.3) {fuente};

  % cascarones
  \draw[figaux] (0,0) circle (1.7);
  \draw[figaux] (0,0) circle (3.2);
  \draw[figvecaux] (0,0) -- (150:1.7);
  \node[figlblaux, anchor=south east] at (150:1.0) {$r_1$};
  \draw[figvecaux] (0,0) -- (170:3.2);
  \node[figlblaux, anchor=south] at (170:2.2) {$r_2$};

  % cono de energía: el mismo haz, dos áreas distintas
  \fill[figamber!14] (0,0) -- (25:1.7) arc (25:55:1.7) -- (55:3.2)
      arc (55:25:3.2) -- cycle;
  \draw[figamber, line width=0.9pt] (25:1.7) arc (25:55:1.7);
  \draw[figamber, line width=0.9pt] (25:3.2) arc (25:55:3.2);
  \draw[figaux] (0,0) -- (25:3.6);
  \draw[figaux] (0,0) -- (55:3.6);
  \node[figlbl, figamber, anchor=north] at (38:1.52) {$A_1$};
  \node[figlbl, figamber, anchor=south] at (38:3.32) {$A_2$};

  % rayos en el resto de las direcciones
  \foreach \ang in {60,80,100,120,250,280,310,340}{
    \draw[figvecaux] (\ang:0.25) -- (\ang:1.45);}

  % lectura
  \node[figlbl, anchor=north west, align=left] at (3.3,-0.5)
    {$P$ es la misma en los dos cascarones:\\[2pt]
     $I_1\,4\pi r_1^2 = I_2\,4\pi r_2^2$\\[2pt]
     $\Rightarrow\; I \propto \dfrac{1}{r^2}, \qquad
      |\mathbf{E}| \propto \dfrac{1}{r}$};

\end{tikzpicture}\\[0.5em]
{\small\itshape La onda plana es cómoda pero imposible: uniforme en todo un
plano infinito, tendría energía total infinita. Lo realista es una fuente
puntual que emite en todas direcciones, y ahí el decaimiento no es una pérdida
de energía sino puro reparto geométrico. La onda no disipa nada al propagarse,
así que la potencia que atraviesa el cascarón chico es la misma que atraviesa el
grande; como la superficie de una esfera crece con $r^2$, la energía por unidad
de área cae con $1/r^2$. El segundo resultado se obtiene enganchando con lo
anterior: lejos de la fuente, en una región chica, la onda esférica se parece
localmente a una plana ---no se alcanza a notar la curvatura--- y ahí vale
$I \propto \langle E^2\rangle$. Si $I$ cae como $1/r^2$, el campo eléctrico cae
como $1/r$, no como $1/r^2$. Vale la pena reparar en la diferencia con el campo
de una carga puntual, que sí va como $1/r^2$: el campo \emph{radiado} decae
mucho más despacio, y por eso la radiación llega lejos.}
\end{center}
